% Generated draft: requires mathematical and source review.
\documentclass[11pt]{article}
\usepackage[margin=25mm]{geometry}
\usepackage{fontspec}
\setmainfont{Latin Modern Roman}
\usepackage{amsmath,amssymb,mathtools,unicode-math}
\setmathfont{Latin Modern Math}
\usepackage{enumitem,needspace,xurl,hyperref}
\setlength{\parindent}{0pt}
\setlength{\parskip}{5pt}
\setlength{\emergencystretch}{3em}
\newcommand{\N}{\mathbb{N}}
\newcommand{\Z}{\mathbb{Z}}
\newcommand{\Q}{\mathbb{Q}}
\newcommand{\R}{\mathbb{R}}
\newcommand{\C}{\mathbb{C}}
\newcommand{\F}{\mathbb{F}}
\newcommand{\A}{\mathbb{A}}
\newcommand{\PP}{\mathbb{P}}
\newcommand{\E}{\mathbb{E}}
\newcommand{\eps}{\varepsilon}
\newcommand{\dd}{\,\mathrm{d}}
\begin{document}
\section*{UnsolvedMath ID: 701173}
\textbf{Source problem code:} catalogue entry
\subsection{Definitions and mathematical statement}
Let \(G\) be a group, with identity element \(1\). It is torsion-free if, for every \(g\in G\) and every integer \(n\geq 1\), the equality \(g^n=1\) implies \(g=1\). For a subgroup \(H\leq G\), define its normalizer by
\[
N_G(H)=\{g\in G:gHg^{-1}=H\}.
\]
The normalizer condition means that
\[
H<N_G(H)\qquad\text{for every proper subgroup }H<G.
\]
Here the condition concerns all proper subgroups, without a finite-generation restriction.

A group \(G\) is hyperabelian if there are an ordinal \(\lambda\) and an ascending series \((A_\alpha)_{\alpha\leq\lambda}\) of normal subgroups of \(G\) such that \(A_0=\{1\}\), \(A_\lambda=G\), every factor \(A_{\alpha+1}/A_\alpha\) is abelian, and
\[
A_\delta=\bigcup_{\alpha<\delta}A_\alpha
\]
for each nonzero limit ordinal \(\delta\leq\lambda\).

Write \(Z(K)\) for the center of a group \(K\). The transfinite upper central series of \(G\) is defined by
\[
Z_0(G)=\{1\},\qquad
Z_{\alpha+1}(G)/Z_\alpha(G)=Z\bigl(G/Z_\alpha(G)\bigr),
\qquad
Z_\delta(G)=\bigcup_{\alpha<\delta}Z_\alpha(G)
\]
for nonzero limit ordinals \(\delta\). The group is hypercentral if \(Z_\lambda(G)=G\) for some ordinal \(\lambda\).

The two questions are:
\begin{enumerate}
\item[(a)] Is every torsion-free group satisfying the normalizer condition hyperabelian?
\item[(b)] Is every torsion-free group satisfying the normalizer condition hypercentral?
\end{enumerate}
No cardinality or finite-generation hypothesis on \(G\) is supplied.

As background, the supplied record attributes to H. Smith the assertion that a torsion-free group all of whose subgroups are subnormal is nilpotent. Here \(H\leq G\) is subnormal if there is a finite chain \(H=H_0\trianglelefteq H_1\trianglelefteq\cdots\trianglelefteq H_m=G\), and nilpotence means that \(Z_c(G)=G\) for some finite integer \(c\geq0\). This background assertion is supplied context, not an additional hypothesis of either question. The registry labels the problem open; that label and the background citation have not been independently verified. The proposer attribution is malformed in the supplied record, as recorded in Sources.
\subsection{Short English statement}
Must a torsion-free group in which every proper subgroup is strictly contained in its normalizer be (a) hyperabelian and (b) hypercentral? These are the two questions in the supplied record, which labels the problem open. This statement-only draft requires human mathematical review.
\subsection{Sources}
\begin{itemize}
\item Supplied local UnsolvedMath registry record, permanent ID \(701173\), exact source code \texttt{catalogue entry}, title \emph{Kourovka Notebook Problem 15.85}. The supplied status is \texttt{open}; the proposed year is 2002. The record was created and updated at \texttt{2026-10-16T14:00:00Z}, with upstream revision \texttt{d371ee917ac9647fbb06cd6fc529deec6f3dda82}.
\item \emph{The Kourovka Notebook (Unsolved Problems in Group Theory)}, problem 15.85, originating in the 15th issue (2002), as identified in the supplied record with the 21st issue edition (2026), \texttt{arXiv:1401.0300v48}. The supplied collection description names E. I. Khukhro and V. D. Mazurov as editors. These bibliographic details are supplied attribution and have not been independently checked.
\item The proposer field is supplied literally as \texttt{Yu. V. Sosnovski\textbackslash{}u\{\textbackslash{}i}. Its trailing character encoding is incomplete; no corrected spelling is inferred.
\item H. Smith, \emph{Arch. Math.}, volume 76, no. 1 (2001), pp. 1--6, cited in the supplied statement for the background assertion about torsion-free groups whose subgroups are all subnormal. No article title or link is supplied.
\end{itemize}
\end{document}
